% concatenated from 260213.tex
\documentclass[11pt]{article}
\usepackage[ngerman, croatian, english]{babel}
\usepackage[T1]{fontenc}
\usepackage[utf8]{inputenc}
\usepackage{geometry}
\geometry{verbose,tmargin=2cm,bmargin=2.5cm,lmargin=3cm,rmargin=3cm}
\usepackage{color}
\usepackage{babel}
\usepackage{textcomp}
\usepackage{mathrsfs}
\usepackage{amsmath}
\usepackage{amsthm}
\usepackage{amssymb}
\usepackage{cancel}
\usepackage{stmaryrd}
\usepackage{stackrel}
\usepackage[all]{xy}
\usepackage[unicode=true,pdfusetitle,
bookmarks=true,bookmarksnumbered=false,bookmarksopen=false,
breaklinks=false,pdfborder={0 0 1},backref=false,colorlinks=true]
{hyperref}
\hypersetup{
	pdfborderstyle=,linkcolor=MidnightBlue, citecolor=PineGreen, urlcolor=PineGreen}

\makeatletter
%%%%%%%%%%%%%%%%%%%%%%%%%%%%%% Textclass specific LaTeX commands.
\theoremstyle{plain}
\newtheorem{thm}{\protect\theoremname}[section]
\theoremstyle{definition}
\newtheorem{defn}[thm]{\protect\definitionname}
\theoremstyle{remark}
\newtheorem{rem}[thm]{\protect\remarkname}
\theoremstyle{plain}
\newtheorem{lem}[thm]{\protect\lemmaname}
\theoremstyle{plain}
\newtheorem{prop}[thm]{\protect\propositionname}
\theoremstyle{plain}
\newtheorem{cor}[thm]{\protect\corollaryname}
\theoremstyle{definition}
\newtheorem*{problem*}{\protect\problemname}
\newtheorem*{funding*}{\protect\fundingname}

\@ifundefined{date}{}{\date{}}
%%%%%%%%%%%%%%%%%%%%%%%%%%%%%% User specified LaTeX commands.
\usepackage[all]{xy}
\usepackage[dvipsnames]{xcolor}
\usepackage{tikz}
\usepackage{hyperref}

\usepackage{csquotes}
\usepackage[backend=biber, style=alphabetic, maxnames=5]{biblatex}
\DeclareFieldFormat{urldate}{}
\DeclareFieldFormat{language}{}
\DeclareFieldFormat{url}{}
\addbibresource{Arrangement.bib}


\usepackage{changepage}

\usepackage{graphicx}
\usepackage{xstring}
\usepackage{forloop}
\usepackage{bbm, amsmath, amsfonts, amscd, latexsym, amsthm, amssymb, graphicx, mathrsfs, comment}
\usepackage{mathtools}
\usepackage{caption}
\usepackage{MnSymbol}
\usepackage{booktabs}


\renewcommand\texteuro{FIXME}

\makeatother

\providecommand{\corollaryname}{Corollary}
\providecommand{\definitionname}{Definition}
\providecommand{\lemmaname}{Lemma}
\providecommand{\problemname}{Problem}
\providecommand{\propositionname}{Proposition}
\providecommand{\remarkname}{Remark}
\providecommand{\theoremname}{Theorem}
\providecommand{\fundingname}{Funding}

\begin{document}
\title{Combinatorics behind discriminants of polynomial systems}
\author{Vladislav Pokidkin}
\maketitle
\begin{abstract}
We develop certain combinatorial tools for the study of discriminants
of general systems of polynomial equations. Applying these tools in
a sequel paper, we completely classify components of such discriminants,
generalizing the classical results of Gelfand, Kapranov, and Zelevinsky
on discriminants of one general multivariate polynomial.

The developed tools are targeted at vector subspace arrangements and
naturally extend to their combinatorial abstraction called polymatroids,
which are the subject matter of this work.

We explore relations between polymatroids and their induced matroids
for bases, circuits, cycles, and rank functions. We define contractions
for polymatroids corresponding to contractions of the induced matroids.
With a view towards applications to discriminants, we construct a
new combinatorial structure induced by polymatroids, called BK-sets.
\end{abstract}
\textbf{Keywords:} polymatroids, induced matroids, discriminants of
polynomial systems.

\section*{Introduction}

Combinatorics of tuples of Newton polytopes is displayed in the structure
of solutions for generic polynomial systems \cite{bernshtein_number_1975,khovanskii_newton_1978,esterov_multivariate_2016,khovanskii_newton_2016},
the intersection theory \cite{fulton_introduction_1993,kaveh_mixed_2010},
the theory of resultants \cite{sturmfels_newton_1994,jensen_computing_2013},
and discriminants \cite{gelfand_discriminants_1994,esterov_newton_2010,cattani_mixed_2013,esterov_discriminant_2013,esterov_galois_2019,dickenstein_iterated_2023}.
The current paper develops such combinatorics and provides a sufficient
background to enumerate components of discriminants of polynomial
systems for the sequel work \cite{pokidkin_components_2025}.

The minimal vector subspace to which a given polytope can be shifted
is said to be generated by this polytope. In that sense, a tuple of
Newton polytopes generates a tuple of real subspaces. The tuple of
vector subspaces can be considered as a realizable polymatroid. It
turned out that the necessary combinatorial results for discriminants
hold for non-realizable polymatroids as well.

In 1970, Endomds introduced the notion of polymatroids. An (integer)
\textit{polymatroid} $P$ is the pair $(E,f)$ of a finite \textit{ground
set} $E$ and a rank function $f:2^{E}\rightarrow\mathbb{Z}_{\geq0}$,
which is submodular, monotone, and normalized ($f(\varnothing)=0$).
The theory of polymatroids is developing \cite{csirmaz_cyclic_2020,crowley_bergman_2022,bonin_natural_2023,swartz_polymatroids_2024},
and it interlinks with algebraic geometry \cite{pagaria_hodge_2023,eur_intersection_2024,baker_representation_2025}.

The \textit{defect} of a subset $I$ from the ground set $E$ is the
difference $\delta(I)=f(I)-|I|$.

In a polymatroid, a set is \textit{essential} if every proper subset
has a strictly greater defect than the ground set. A set is\textit{
independent} if the defects of all subsets are non-negative. An independent
set is \textit{irreducible} if the defects of all proper subsets are
positive. A \textit{BK-set} is an independent set with zero defect.

We characterize polymatroids via their ground sets. For example a
dependent polymatroid is a polymatroid with a dependent ground set,
and a BK-polymatroid is a polymatroid for which the ground set is
a BK-set.

Every polymatroid $P=(E,f)$ induces the matroid $M=(E,rk)$ \cite{mcdiarmid_independence_1973,welsh_matroid_1986}
such that 
\[
rk(I)=min\,\{f(J)+|I\backslash J|,\;J\subseteq E\}.
\]

The induced matroid is a contraction of the natural matroid of a polymatroid
\cite{bonin_natural_2023}. 

The main results are:
\begin{enumerate}
\item General results about induced matroids:
\begin{enumerate}
\item (In)dependent sets in a polymatroid are (in)dependent sets in the
induced matroid (Theorem \ref{Theorem. Independent in polymatroid iff idependent in matroid}
and Proposition \ref{Proposition. Dependent in polymatroid dependent in matroid}).
\item All circuits have the defect $-1$ (Corollary \ref{Corollary. Defect of circuits is -1}).
\item All bases of the induced matroid have their unique maximal BK-sets
of the same cardinality for a given polymatroid (Theorem \ref{Theorem. Each basis contains a unique maxamal BK-subtuple}
and Corollary \ref{Corollary. BK-sets in bases have the same cardinality}).
\end{enumerate}
\item Results about cycles:
\begin{enumerate}
\item Bases are BK-sets (Theorem \ref{Theorem. In cyclic, bases are BK}).
\item A subset is cyclic if and only if it is essential (Theorem \ref{Theorem. Cyclic=00003Dessential}).
\item The rank of the induced matroid equals the polymatroid rank (Corollary
\ref{Corollary. cyclic, rk(I)=00003Df(I)}).
\item There exists a unique subset such that the following are equivalent
(Theorem \ref{Theorem. Maximal essential =00003D minimal defect}):
\begin{enumerate}
\item the subset is maximal essential,
\item the subset is maximal cyclic,
\item the subset has a minimal defect and is minimal by inclusion.
\end{enumerate}
\end{enumerate}
\item Results about ranks:
\begin{enumerate}
\item A new formula for the rank of the induced matroid is obtained. This
formula explicitly indicates which cyclic flat we need to use in the
Sims, Bonin, and Mier formula \cite{sims_problems_1980,bonin_lattice_2008}
(Proposition \ref{Proposition. New formulat for rank}). 
\item All subsets containing a basis of the induced matroid have the same
polymatroid rank (Theorem \ref{Theorem. f(B)=00003Df(E)}).
\item The polymatroid rank equals the induced matroid rank for the ground
set if and only if the bases of the induced matroid are BK-sets (Corollary
\ref{Corollary. rk(E)=00003Df(E) iff bases are BK}).
\end{enumerate}
\item Results about contractions:
\begin{enumerate}
\item Contractions of polymatroids are defined, and they correspond to the
contractions of the induced matroids:
\begin{enumerate}
\item Contraction by a BK-set for an independent polymatroid (Proposition
\ref{Proposition. Polymatroid contraction by BK});
\item Contraction by a cyclic flat for a dependent polymatroid (Theorem
\ref{Theorem. Contraction by a cyclic flat}).
\end{enumerate}
\item The contraction of a polymatroid by its maximal cycle is independent
(Corollary \ref{Corollary. Contraction by maximal cyclic is independent}).
\end{enumerate}
\item Results about BK-sets:
\begin{enumerate}
\item For a BK-set in the induced matroid, the set of BK-subsets forms a
distributive lattice by inclusion (Corollary \ref{Corollary. Distributice lattice of BK-subsets}).
Therefore, BK-sets are endowed with the antimatroid structure \cite{edelman_meet-distributive_1980}.
\item Every BK-set admits a partition, encoded by the Birkhoff poset, which
can be arbitrary (Theorem \ref{Theorem. The poset can be arbitratry}).
Every subset of the partition corresponds to an irreducible BK-polymatroid
(Theorem \ref{Theorem. BK-partition}).
\end{enumerate}
\end{enumerate}
With a view towards applications to discriminants, we construct a
certain partition of a vector space over an infinite field by a realizable
polymatroid (Proposition \ref{Proposition. Polymatroid partition of a vector space}).
Section \ref{Section. Disctriminants} discusses the meaning of results
in the context of discriminants of polynomial systems.

\section{Polymatroids and induced matroids}

\subsection{Independent sets in polymatroids and induced matroids}
\begin{thm}
\label{Theorem. Independent in polymatroid iff idependent in matroid}A
polymatroid is independent if and only if the induced matroid is a
boolean algebra.
\end{thm}

\begin{proof}
$\boxed{\Rightarrow}$ It is enough to show that $rk(I)=|I|$ for
every subset $I\subseteq E$. Notice that the inequality $rk(I)\leq|I|$
holds automatically. By contradiction, suppose the strict inequality
$rk(I)<|I|.$ Then there exists a subset $J\subseteq E$ reaching
the rank for the induced matroid
\[
rk(I)=f(J)+|I\backslash J|\geq|J|+|I\backslash J|=|I\cup J|\geq|I|.
\]
The first inequality holds because the polymatroid is independent,
and this leads to the contradiction.

$\boxed{\Leftarrow}$ If the induced matroid is a boolean algebra,
then $rk(I)=|I|$ for every subset $I\subseteq E$: 
\[
f(I)\geq min\,\{f(J)+|I\backslash J|,\;J\subseteq E\}=|I|.
\]
This means that the defects of all subsets are non-negative.
\end{proof}
\begin{cor}
An independent subset in a polymatroid is an independent subset in
the induced matroid.
\end{cor}

\begin{prop}
\label{Proposition. Dependent in polymatroid dependent in matroid}A
set is dependent in a polymatroid if and only if the set is dependent
in the induced matroid.
\end{prop}

\begin{proof}
$\boxed{\Rightarrow}$ A set is dependent in a polymatroid if it contains
a subset $I$ of a negative defect, $f(I)-|I|<0$. Since the polymatroid
rank is not smaller than the induced matroid rank $f(I)\geq rk(I)$,
we can conclude $rk(I)-|I|<0,$ and the set $I$ is dependent in the
induced matroid.

$\boxed{\Leftarrow}$ Consider a dependent set $I$ in the induced
matroid, $rk(I)-|I|<0$. If we restrict the induced matroid to the
set $I$, the rank of $I$ does not change. The inequality is equivalent
to 
\[
rk(I)=min\,\{f(J)+|I\backslash J|,\;J\subseteq I\}=f(S)+|I\backslash S|<|I|
\]
for some choice of a subset $S\subseteq I$. Therefore, the set $I$
contains the subset $S$ of a negative defect, and the set $I$ is
dependent in the polymatroid.
\end{proof}
\begin{cor}
In a polymatroid, subsets of negative defects correspond to dependent
subsets in the induced matroid.
\end{cor}


\subsection{Defects of circuits for induced matroids}
\begin{lem}
\label{Lemma. In essential, f(E)=00003Df(E=00005Ca)}For an essential
polymatroid and an element $e$ from the ground set $E$, it holds:
$f(E)=f(E\backslash e).$
\end{lem}

\begin{proof}
Since the ground set is essential, $\delta(E\backslash e)>\delta(E)$
by definition. This is equivalent to $1>f(E)-f(E\backslash e)$. Since
the function $f$ is monotone, we can have only the equality.
\end{proof}
\begin{lem}
\label{Lemma. Circuits have negative defect}Circuits of the induced
matroid from a polymatroid have negative defects.
\end{lem}

\begin{proof}
For a circuit, every proper subset is independent and has a non-negative
defect. By Proposition \ref{Proposition. Dependent in polymatroid dependent in matroid},
a dependent subset contains a subset of a negative defect. Hence there
is only one option for circuits: the defect of the whole circuit is
negative.
\end{proof}
\begin{lem}
\label{Lemma. C=00005Ce is a BK-set}For an element $e$ from a circuit
$C$ of the induced matroid from a polymatroid, the complement $C\backslash e$
is a BK-set.
\end{lem}

\begin{proof}
It is sufficient to show that the independent set $C\backslash e$
has zero defect. Theorem \ref{Theorem. Independent in polymatroid iff idependent in matroid}
implies that circuits are essential sets. Use Lemma \ref{Lemma. In essential, f(E)=00003Df(E=00005Ca)},
write the inequalities
\[
|C\backslash e|\overset{C\backslash e\text{ is independent}}{\leq}f(C\backslash e)=f(C)\overset{C\text{ is dependent}}{<}|C|,
\]
and conclude that the defect of $C\backslash e$ is zero.
\end{proof}
\begin{cor}
\label{Corollary. Defect of circuits is -1}Circuits of the induced
matroid have the defect $-1$.
\end{cor}


\subsection{Defects of bases for induced matroids}
\begin{prop}
\textup{\label{Proposition. Oxley Circuits}\cite{oxley_matroid_2011}}
Suppose $I$ is an independent set in a matroid and $e$ is an element
such that $I\cup e$ is dependent; then there exists a unique circuit
contained in $I\cup e$, and this circuit contains $e$.
\end{prop}

\begin{lem}
\label{Lemma. Value of f equal}Suppose $I$ is an independent set
in the induced matroid from a polymatroid and $e$ is an element such
that $I\cup e$ is dependent; then the values are equal: $f(I)=f(I\cup e).$
\end{lem}

\begin{proof}
By Proposition \ref{Proposition. Oxley Circuits}, there exists a
circuit $C\subseteq I\cup e$ such that it contains the element $e$.
By Lemma \ref{Lemma. In essential, f(E)=00003Df(E=00005Ca)}, $f(C)=f(C\backslash e).$
The polymatroid rank $f$ is submodular, and
\begin{align*}
f(C)+f(I) & \geq f(I\cup C)+f(I\cap C)=f(I\cup e)+f(C\backslash e),\\
f(I) & \geq f(I\cup e).
\end{align*}
Since the polymatroid rank $f$ is monotone, we have the equality
$f(I)=f(I\cup e).$
\end{proof}
Recall the basis-exchange property for matroids: if bases $B$ and
$B'$ are distinct, and $e\in B\backslash B'$, then there exists
an element $e'\in B'\backslash B$ such that $(B\cup e')\backslash e$
is a basis.
\begin{lem}
\label{Lemma. Basis change}Basis change does not change the defect:
$\delta(B)=\delta((B\cup e')\backslash e).$
\end{lem}

\begin{proof}
Use Lemma \ref{Lemma. Value of f equal} twice:

$\delta((B\cup e')\backslash e)=f((B\cup e')\backslash e)-|(B\cup e')\backslash e|=f(B\cup e')-|B|=f(B)-|B|=\delta(B).$
\end{proof}
\begin{thm}
\label{Theorem. Basis have the same defect}Bases of the induced matroid
have the same non-negative defect.
\end{thm}

\begin{proof}
By Lemma \ref{Lemma. Basis change}, the basis exchange property preserves
the defect. Since bases are linked with each other by a sequence of
basis changes, all bases have the same defect. The defect is non-negative,
because bases are independent.
\end{proof}
\begin{cor}
Bases of the induced matroid have the same polymatroid rank.
\end{cor}


\subsection{BK-sets in a basis}
\begin{lem}
\label{Lemma. The union and intersection of BK-sets are BK-sets}In
an independent polymatroid, unions and intersections of BK-sets are
BK-sets.
\end{lem}

\begin{proof}
In an independent polymatroid, the defects of all subsets are non-negative.
The polymatroid rank function $f$ is submodular:
\begin{align*}
f(I\cup J)+f(I\cap J) & \leq f(I)+f(J)\overset{I,J\text{ are BK}}{=}|I|+|J|=|I\cup J|+|I\cap J|\;\Leftrightarrow\\
\delta(I\cup J)+\delta(I\cap J) & \leq0.
\end{align*}
For non-negative defects, there is only one option: $\delta(I\cup J)=\delta(I\cap J)=0.$
\end{proof}
\begin{cor}
\label{Corollary. Distributice lattice of BK-subsets}For an independent
polymatroid, BK-sets form a distributive lattice by inclusion.
\end{cor}

\begin{rem}
Edelman established that distributive lattices are antimatroids \cite{edelman_meet-distributive_1980}.
This means that each basis of the induced matroid from a polymatroid
is equipped with the antimatroid structure.
\end{rem}

\begin{lem}
\label{Lemma. Ecah basis contains a BK-set}For a dependent polyamtroid,
every basis of the induced matroid contains a non-trivial BK-set.
\end{lem}

\begin{proof}
For a basis $B$ and an element $e$ from the complement $E\backslash B$,
there exists a unique circuit $C$ in $B\cup e$ containing $e$ according
to Proposition \ref{Proposition. Oxley Circuits}. By Lemma \ref{Lemma. C=00005Ce is a BK-set},
the subset $C\backslash e$ is a BK-set in the basis $B$.
\end{proof}
\begin{thm}
\label{Theorem. Each basis contains a unique maxamal BK-subtuple}For
a dependent polymatroid, every basis of the induced matroid contains
a unique BK-set maximal by inclusion.
\end{thm}

\begin{proof}
By Lemma \ref{Lemma. Ecah basis contains a BK-set}, every basis contains
a non-trivial BK-set. According to Lemma \ref{Lemma. The union and intersection of BK-sets are BK-sets},
the union of all BK-sets in a basis is a unique and maximal BK-set
by inclusion.
\end{proof}

\section{Cycles of the induced matroid}

\subsection{Bases of cycles}
\begin{lem}
\label{Lemma. Union with a circuit decreases the defect}For a circuit
$C$ that does not lie in a subset $I$, $\delta(I\cup C)<\delta(C).$
\end{lem}

\begin{proof}
For such a circuit $C$, the intersection $I\cap C$ is independent
of a non-negative defect.
\begin{align*}
\delta(I\cup C) & =f(I\cup C)-|I\cup C|=\\
 & =\delta(I)+\underbrace{\delta(C)}_{-1}-\underbrace{\delta(I\cap C)}_{\geq0}-\underbrace{(f(I)+f(C)-f(I\cup C)-f(I\cap C))}_{\geq0}<\delta(I)
\end{align*}
We use Corollary \ref{Corollary. Defect of circuits is -1} and the
submodularity of $f$.
\end{proof}
\begin{cor}
A union of circuits has a negative defect.
\end{cor}

Recall that a \textit{cycle} \textit{(cyclic} \textit{subset)} is
a union of circuits, and a \textit{loop}/\textit{coloop} is an element
contained in no/every basis.
\begin{cor}
\label{Corollary. Cyclics are essential}In the induced matroid, cycles
are essential.
\end{cor}

\begin{lem}
\label{Lemma. Cyclics are unions of BK-sets}In the induced matroid,
cycles are unions of BK-sets.
\end{lem}

\begin{proof}
By Lemma \ref{Lemma. C=00005Ce is a BK-set}, every circuit $C$ is
a union of $|C|$-number BK-sets.
\end{proof}
\begin{thm}
\label{Theorem. In cyclic, bases are BK}In a cyclic polymatroid,
bases are BK-sets.
\end{thm}

\begin{proof}
A cyclic ground set $E$ is a union of BK-sets by Lemma \ref{Lemma. Cyclics are unions of BK-sets}
and is a union of circuits by definition. For a fixed element $e\in E,$
we can represent the ground set $E$ via the union of circuits as
follows:
\[
E=\left(\underset{e\in C}{\cup}C\right)\cup\left(\underset{e\notin C'}{\cup}C'\right).
\]
The union $\underset{e\notin C'}{\cup}C'$ is a cycle; hence, it is
a union of BK-sets by Lemma \ref{Lemma. Cyclics are unions of BK-sets}.
According to Lemma \ref{Lemma. C=00005Ce is a BK-set}, $C\backslash e$
is a BK-set for every circuit $C$ containing $e.$ Therefore, the
complement
\[
E\backslash e=\left(\underset{e\in C}{\cup}(C\backslash e)\right)\cup\left(\underset{e\notin C'}{\cup}C'\right)
\]
is a union of BK-sets. If the complement $E\backslash e$ is dependent,
we repeat the procedure and remove one by one elements until we obtain
an independent set $I\subset E$. This independent set $I$ is a basis
automatically, is a union of BK-sets by induction, and is a BK-set
by Lemma \ref{Lemma. The union and intersection of BK-sets are BK-sets}.
\end{proof}
\begin{cor}
If bases of the induced matroid from a polymatroid have a positive
defect, then the induced matroid contains a coloop.
\end{cor}

\begin{cor}
\label{Corollary. BK-sets in bases have the same cardinality}In a
dependent polymatroid, the unique maximal BK-sets in the bases of
the induced matroid have the same cardinality.
\end{cor}

\begin{proof}
For the maximal cycle $E'$ in the ground set $E$, the complement
$E\backslash E'$ consists of coloops. Then every basis $B$ of the
induced matroid is the union $B'\cup(E\backslash E'),$ where $B'=B\cap E'$.
The set $B'$ is a basis for the restricted matroid on $E'$; hence,
$B'$ is a BK-set by Theorem \ref{Theorem. In cyclic, bases are BK}.
Thus the set $B'$ is contained in the maximal BK-set $I$ in the
basis $B$, and $B'=I\cap E'$. This means that the complement $J=I\backslash B'$
consists of coloops. Then the set $J$ is contained in every basis
of the induced matroid on $E$. For the restriction of the induced
matroid to the set $E'\cup J$, the set $I$ is a basis and a BK-set.
Theorem \ref{Theorem. Basis have the same defect} implies that all
bases for the set $E'\cup J$ are BK-sets. Therefore, for every basis
$B$ on $E$, its maximal BK-set $I$ has the same fixed cardinality,
and it equals the union $B'\cup J,$ where $B'=B\cap E'$ is a basis
of $E'$.
\end{proof}

\subsection{Cyclic is equivalent to essential}
\begin{prop}
\label{Proposition. Essential, rk(E)=00003Df(E)}For an essential
polymatroid, the rank of the induced matroid of the ground set $E$
equals the polymatroid rank, $rk(E)=f(E).$
\end{prop}

\begin{proof}
For an essential ground set, the defects of proper subsets are strictly
greater than the defect of the ground set. Then
\[
rk(E)=min\,\{f(J)+|E\backslash J|,\;J\subseteq E\}=|E|+min\,\{\delta(J),\;J\subseteq E\}=|E|+\delta(E)=f(E).
\]
\end{proof}
\begin{prop}
\textup{\label{Proposition. Characterization of coloops}\cite{white_theory_1986}
}An element $e$ from the ground set $E$ of a matroid is a coloop
if and only if it satisfies one of the following:

Bases: $e$ is in every basis;

Circuits: $e$ is in no circuit;

Rank: $rk(E\backslash e)=rk(E)-1$.
\end{prop}

\begin{lem}
\label{Lemma. An essential dependent polymatroid doesn't contain coloops.}For
an essential dependent polymatroid, the induced matroid does not contain
coloops.
\end{lem}

\begin{proof}
For the maximal cycle $E'$ in the essential ground set $E$, the
complement $E\backslash E'$ consists of coloops. By Proposition \ref{Proposition. Characterization of coloops},
we can write the equality for ranks: $rk(E)-rk(E')=|E\backslash E'|$.
Since the set $E'$ is a cycle, the bases of $E'$ are BK-set by Theorem
\ref{Theorem. In cyclic, bases are BK}. Notice that the matroid rank
does not exceed the polymatroid rank, $rk(E)\leq f(E)$, use Corollary
\ref{Corollary. Cyclics are essential} and Proposition \ref{Proposition. Essential, rk(E)=00003Df(E)},
$rk(E')=f(E')$, and conclude
\[
f(E)-f(E')\geq|E\backslash E'|.
\]
This inequality is equivalent to $\delta(E)\geq\delta(E')$. Therefore,
the ground set is not essential if it contains coloops.
\end{proof}
\begin{thm}
\label{Theorem. Cyclic=00003Dessential}A dependent polymatroid is
essential if and only if it is cyclic.
\end{thm}

\begin{proof}
Corollary \ref{Corollary. Cyclics are essential} and Lemma \ref{Lemma. An essential dependent polymatroid doesn't contain coloops.}.
\end{proof}
\begin{cor}
\label{Corollary. cyclic, rk(I)=00003Df(I)}For a dependent polymatroid,
the rank of the induced matroid of a cycle $I$ equals the polymatroid
rank, $rk(I)=f(I).$
\end{cor}

\begin{proof}
Restrict the induced matroid to the set $I$ and use Proposition \ref{Proposition. Essential, rk(E)=00003Df(E)}
and Theorem \ref{Theorem. Cyclic=00003Dessential}.
\end{proof}
\begin{defn}
The \textit{matroid} $cl_{M}(I)$/\textit{polymatroid} $cl_{P}(I)$
\textit{closure} of a set $I$ is a maximal subset containing $I$
of the same matroid/polymatroid rank.
\end{defn}

\begin{cor}
\label{Corollary. Cyclic, two closures are equal}For a cycle $I$
from a dependent polymatroid, the polymatroid closure $cl_{P}(I)$
equals the induced matroid closure $cl_{M}(I)$.
\end{cor}

\begin{proof}
By Corollary \ref{Corollary. cyclic, rk(I)=00003Df(I)}, the rank
functions are equal in that case.
\end{proof}

\subsection{Maximal essential subset}
\begin{prop}
\label{Proposition. Unique minimal subset of the minimal defect}A
dependent polymatroid contains the unique subset with minimal defect
and minimal by inclusion.
\end{prop}

\begin{proof}
Suppose there exist two distinct such sets, $I$ and $J$; then the
defect of the intersection $I\cap J$ is strictly greater than the
defects of $I$ and $J$. We obtain a contradiction, since the union
$I\cup J$ has a strictly smaller defect:
\begin{align*}
\delta(I\cup J) & =\delta(I)+\underbrace{\delta(J)-\delta(I\cap J)}_{<0}-\underbrace{(f(I)+f(J)-f(I\cup J)-f(I\cap J))}_{\geq0}<\delta(I)=\delta(J).
\end{align*}
\end{proof}
\begin{thm}
\label{Theorem. Maximal essential =00003D minimal defect}In a dependent
polymatroid, the following are equivalent:

a) the subset is maximal essential,

b) the subset is maximal cyclic in the induced matroid,

c) the subset has a minimal defect and minimal by inclusion.
\end{thm}

\begin{proof}
$a)=b)$ by Theorem \ref{Theorem. Cyclic=00003Dessential}. $b)\subseteq c)$
by Lemma \ref{Lemma. Union with a circuit decreases the defect}.

$c)\Rightarrow a)$ by Proposition \ref{Proposition. Unique minimal subset of the minimal defect},
the subset $I$ of minimal defect and minimal by inclusion is unique.
Hence every subset containing $I$ can not be essential by definition.
\end{proof}

\section{Polymatroid and the induced matroid ranks}

Bonin and Mier showed that the set of cyclic flats forms an arbitrary
lattice \cite{bonin_lattice_2008}. The join $\vee$ of cyclic flats
is a closure of the union, and the meet $\wedge$ of cyclic flats
is a maximal cyclic flat inside the intersection. The lattice of cyclic
flats admits an alternative cryptomorphic definition for matroids.

\subsection{Rank of the induced matroid}
\begin{thm}
\textup{\label{Theorem. Sims, Bonin, Mier, rank funciton}(Sims \cite{sims_problems_1980},
Bonin, Mier} \textup{\cite{bonin_lattice_2008})} For a subset $I$
from a matroid with the lattice of cyclic flats $\mathcal{Z}$, the
matroid rank function satisfies the equality:
\[
rk(I)=min\,\{rk(F)+|I\backslash F|,\;F\in\mathcal{Z}\}.
\]
\end{thm}

\begin{prop}
\label{Proposition. New formulat for rank}The rank function of the
induced matroid and the polymatroid rank function are linked with
each other as follows:
\[
rk(I)=f(I')+|I\backslash I'|=|I|+\delta(I'),
\]
where $I'$ is the maximal cycle inside $I$.
\end{prop}

\begin{proof}
If we restrict the induced matroid to the set $I$, the rank does
not change and equals 
\[
rk(I)=min\,\{f(J)+|I\backslash J|,\;J\subseteq I\}=|I|+min\,\{\delta(J),\;J\subseteq I\}=|I|+\delta(I'),
\]
where $I'$ is the maximal cycle inside $I$, which is automatically
flat. The last equality holds by Theorem \ref{Theorem. Maximal essential =00003D minimal defect}.
\end{proof}
\begin{rem}
Proposition \ref{Proposition. New formulat for rank} is coherent
with the formula of Sims, Bonin, and Mier. Firstly, for the lattice
of cyclic flats $\mathcal{Z}$ and a subset $I$ with its maximal
cycle $I'$, we need to choose the matroid closure $cl_{M}(I')$ to
obtain an element from the lattice $\mathcal{Z}$. Secondly, the rank
of the induced matroid equals the polymatroid rank for cycles by Corollary
\ref{Corollary. cyclic, rk(I)=00003Df(I)}, $rk(I')=f(I')$. By the
same reason, thirdly, the polymatroid closure of $cl_{P}(I')$ equals
the matroid closure of $cl_{M}(I')$, and 
\[
rk(I')=rk(cl_{M}(I'))=f(cl_{M}(I'))=f(cl_{P}(I'))=f(I').
\]
\end{rem}


\subsection{Polymatroid rank plateau}
\begin{thm}
\label{Theorem. f(B)=00003Df(E)}All subsets containing a basis of
the induced matroid from a polymatroid have the same polymatroid rank.
In particular, $f(B)=f(E)$ for a basis $B$ from the ground set $E$.
\end{thm}

\begin{proof}
Since the polymatroid rank $f$ is monotone, it is enough to prove
that $f(B)=f(E)$. For the maximal cycle $E'$ in the ground set $E$,
the complement $E\backslash E'$ consists of coloops. Choose a basis
$B$ and denote by $B'$ the intersection $B\cap E'$. Notice that
the set of coloops is the same $E\backslash E'=B\backslash B'$, and
we have the equality for the union: $E=B\cup E'$. Notice that $B'$
is a basis in the essential set $E'$, and
\[
f(E')=rk(E')=rk(B')=|B'|=f(B')
\]
by Proposition \ref{Proposition. Essential, rk(E)=00003Df(E)}, the
definition of a matroid rank, and Theorem \ref{Theorem. In cyclic, bases are BK}
correspondingly. Since the rank $f$ is submodular
\[
f(E')+f(B)\geq f(E)+f(B'),
\]
we obtain $f(B)\geq f(E)$. Since $f$ is monotone, we obtain the
equality: $f(B)=f(E).$
\end{proof}
\begin{cor}
\label{Corollary. rk(E)=00003Df(E) iff bases are BK}The polymatroid
rank of the ground set equals the induced matroid rank if and only
if bases are BK-sets.
\end{cor}

\begin{proof}
By Theorem \ref{Theorem. Basis have the same defect}, bases have
the same defect. Choose a basis $B$ and use Theorem \ref{Theorem. f(B)=00003Df(E)}.

$\boxed{\Rightarrow}$ $|B|=rk(E)=f(E)=f(B)$. $\boxed{\Leftarrow}$
$rk(E)=|B|=f(B)=f(E)$.
\end{proof}

\section{Polymatroid contraction}

\subsection{Polymatroid contraction by a BK-set}
\begin{defn}
The \textit{contraction} $P/K$ of a polymatroid $P$ by a BK-set
$K$ is a polymatroid $(E\backslash K,f_{K})$ with the rank function
$f_{K}(I)=f(I\cup K)-f(K).$
\end{defn}

We use the index to emphasize defects after a contraction: $\delta_{K}(I)=f_{K}(I)-|I|$.
\begin{lem}
\label{Lemma. Contraction by BK does not change defect}Contraction
by a BK-set $K$ does not change the defect, $\delta(I)=\delta_{K}(I)$.
\end{lem}

\begin{proof}
$\delta_{K}(I)=f_{K}(I)-|I\backslash K|=f(I)-f(K)-|I\backslash K|=\delta(I).$
\end{proof}
\begin{cor}
For the contraction $P/K$ of a polymatroid $P$ by a BK-set $K$,
there is a bijection between BK-sets in the contraction $P/K$ and
BK-sets containing $K$ in the polymatroid $P$.
\end{cor}

\begin{prop}
\label{Proposition. Polymatroid contraction by BK}The contraction
of a polymatroid by a BK-set corresponds to the contraction of the
induced matroid by the BK-set.
\end{prop}

\begin{proof}
$\boxed{\Rightarrow}$ Consider an independent subset $I$ in the
contraction $P/K$. Let us show that the union $I\cup K$ is independent
in $P$. This is equivalent to show that every subset $J\subseteq I\cup K$
has a non-negative defect. Since the set $I$ is independent in $P/K$
and $K$ is independent in $P,$ we have 
\begin{align*}
f_{K}(I\cap J)-|I\cap J| & \geq0,\\
f(K\cap J)-|K\cap J| & \geq0.
\end{align*}
Sum the inequalities and use the definition of a contraction to obtain
\[
f(K\cup J)-f(K)+f(K\cap J)-|J|\geq0.
\]
Recall that the function $f$ is submodular, $f(J)\geq f(K\cup J)-f(K)+f(K\cap J)$,
and conclude that the defect of $J$ is non-negative in $P$. Therefore,
the set is $I\cup K$ is independent in $P$.

$\boxed{\Leftarrow}$ Notice that independent sets in $P$ corresponds
to independent sets in $P/K$, because the contraction by a BK-set
preserves defects by Lemma \ref{Lemma. Contraction by BK does not change defect}.
\end{proof}
\begin{rem}
The contraction of a polymatroid by a BK-set is preferable for independent
polymatroids, because in that case, we do not have elements of zero
value after the contraction.
\end{rem}


\subsection{Polymatroid contraction by a cyclic flat}

To avoid elements of zero polymatroid value after contraction, we
define contractions by cyclic flats.
\begin{defn}
The \textit{contraction $P/F$} of a polymatroid $P$ by a cyclic
flat $F$ is a polymatroid $(E\backslash F,f_{F})$ with the rank
function $f_{F}(I)=f(I\cup F)-f(F).$
\end{defn}

\begin{lem}
\label{Lemma. Defect of a contraction by a cyclic flat}The contraction
$P/F$ by a cyclic flat has the defect $\delta_{F}(I)=\delta(I)-\delta(F)$.
\end{lem}

\begin{lem}
\label{Lemma. Contraction of a cyclic flat is cyclic flat}Polymatroid
contraction by a cyclic flat corresponds to the matroid contraction
by the cyclic flat on the lattice of cyclic flats. In particular,
the polymatroid contraction by a cyclic flat of a cyclic flat is a
cyclic flat.
\end{lem}

\begin{proof}
In matroids, the contraction of a cyclic flat by a cyclic flat is
a cyclic flat. By Corollary \ref{Corollary. cyclic, rk(I)=00003Df(I)},
the polymatroid rank equals the matroid rank on the lattice of cyclic
flats $\mathcal{Z}$ of the induced matroid. Then the two types of
contractions are equal on $\mathcal{Z}$
\[
f_{F}(I)=f(I\vee F)-f(F)=rk(I\vee F)-rk(F)=rk_{F}(I),
\]
where $I$ and $F$ are cyclic flats, and $rk_{F}(I)$ is the rank
of the matroid contraction of the induced matroid by $F$. We have
applied Theorem \ref{Theorem. f(B)=00003Df(E)} to use the equality
$f(I\cup F)=f(I\vee F)$. Hence the polymatroid contraction corresponds
to the matroid contraction on the lattice of cyclic flats $\mathcal{Z}$.
\end{proof}
Notice that the inclusion $I'\cup F'\subseteq(I\cup F)'$ can be strict
(the sign $'$ means the maximal cycle inside a set).
\begin{prop}
\label{Proposition. Formula for the maximal cyclic in a contraction}For
a set $I$ in the contraction $P/F$, its maximal cycle $I_{F}'$
in the contraction $P/F$ can be computed via the maximal cycle $(I\cup F)'$
inside the union $I\cup F$ in the polymatroid $P$:
\[
I'_{F}=(I\cup F)'\backslash F.
\]
\end{prop}

\begin{proof}
$\boxed{\subseteq}$ It follows from the inclusion $I'_{F}\cup F'\subseteq(I\cup F)'$
in the ground set $E$. $\boxed{\supseteq}$ By Lemma \ref{Lemma. Contraction of a cyclic flat is cyclic flat},
the contraction of the cyclic flat $(I\cup F)'$ is the cyclic flat
$(I\cup F)'\backslash F$ that lies inside the maximal cyclic flat
$I'_{F}$ of the set $I$.
\end{proof}
\begin{thm}
\label{Theorem. Contraction by a cyclic flat}The contraction of a
polymatroid by a cyclic flat corresponds to the contraction of the
induced matroid by the cyclic flat.
\end{thm}

\begin{proof}
Consider a polymatroid $P$ and a cyclic flat $F$ from the induced
matroid $M(P).$ It is enough to show that the diagram is commutative:
\[
\xymatrix{P\ar[r]\ar[d] & M(P)\ar[d]\\
P/F\ar[r] & M(P/F).
}
\]
The diagram is commutative if the two types of ranks coincide: $rk_{F}(I)=rk(I\cup F)-rk(F),$
where $F$ is a cyclic flat and $I$ is a set in the complement $E\backslash F$.
Use Corollary \ref{Corollary. cyclic, rk(I)=00003Df(I)}, Proposition
\ref{Proposition. New formulat for rank}, Proposition \ref{Proposition. Formula for the maximal cyclic in a contraction}
and definitions:
\begin{align*}
rk_{F}(I) & =f_{F}(I'_{F})+|I\backslash I'_{F}|=f_{F}((I\cup F)'\backslash F)+|I\backslash\left((I\cup F)'\backslash F\right)|=\\
 & =f((I\cup F)')-f(F)+|(I\cup F)\backslash(I\cup F)'|=rk(I\cup F)-rk(F).
\end{align*}
Notice that $(I\cup F)\backslash(I\cup F)'=I\backslash\left((I\cup F)'\backslash F\right)=I\backslash I'_{F}$,
because we have the inclusion for cyclic flats $F\subseteq(I\cup F)'$.
\end{proof}
\begin{rem}
If we build a contraction by a flat $F$ which is not cyclic, then
the two types of ranks from the theorem above can differ.
\end{rem}

\begin{cor}
\label{Corollary. Contraction by maximal cyclic is independent}For
the maximal cycle $E'$ of the induced matroid from a polymatroid
$P$, the contraction of $P/E'$ is independent, and the induced matroid
$M(P/E')$ is a boolean algebra.
\end{cor}

\begin{proof}
The maximal cycle $E'$ is the maximal cyclic flat. By Theorem \ref{Theorem. Contraction by a cyclic flat}
the contraction is well-defined. By Lemma \ref{Lemma. Defect of a contraction by a cyclic flat},
the defects of all sets in the polymatroid $P/E'$ are positive.
\end{proof}

\section{BK-sets}

\subsection{Birkhoff poset of a BK-set}
\begin{defn}
An \textit{order ideal} of a poset $\mathtt{P}$ is a subposet $\mathtt{I}$
of $\mathtt{P}$ such that if $\beta\in\mathtt{I}$ and $\alpha\leq\beta$,
then $\alpha\in\mathtt{I}$. For an element $\alpha$ from a poset,
the \textit{principal order ideal }$(\alpha)$ is an order ideal of
all elements that are not greater than $\alpha.$
\end{defn}

In 1937, Birkhoff proved \cite{birkhoff_rings_1937} that every finite
distributive lattice is isomorphic to the lattice of order ideals
of some poset $\mathtt{P}.$
\begin{thm}
\label{Theorem. The poset can be arbitratry}Any finite poset $\mathtt{P}$
can be a Birkhoff poset for BK-sets of a realizable BK-polymatroid.
\end{thm}

\begin{proof}
For a field of the characteristic zero, the incidence algebra for
the poset $\mathtt{P}$ over $n$ elements (see \cite{stanley_enumerative_2011})
is isomorphic to the subalgebra of upper matrices $L=\{(m_{ij})\in\mathrm{Mat}_{n,n}\,|\;m_{ij}=0$
if $i\nleq j,\;i,j\in\mathtt{P}\}.$ The algebra of upper matrices
admits the decomposition on columns as a vector space $L=\underset{j\in P}{\oplus}L_{j},$
where $L_{j}=\langle m_{ij}\,|\;i\leq j\rangle.$ This decomposition
provides a tuple of vector subspaces and a realiztion of a BK-polymatroid
$P$. By the construction, every order ideal $\mathtt{I}\subseteq\mathtt{P}$
corresponds to a BK-set, and vice versa. Hence the lattice of order
ideals of $\mathtt{P}$ is isomorphic to the lattice of BK-sets of
the polymatroid $P$.
\end{proof}

\subsection{Poset partition of a reducible BK-polymatroid}
\begin{prop}
In an independent polymatroid, BK-sets satisfy the modular law:

$f(I)+f(J)=f(I\cup J)+f(I\cap J)$.
\end{prop}

\begin{proof}
By Lemma $\ref{Lemma. The union and intersection of BK-sets are BK-sets}$,
unions and intersections of BK-sets $I$ and $J$ are BK-sets:

$f(I)+f(J)-f(I\cup J)-f(I\cap J)=\delta(I)+\delta(J)-\delta(I\cup J)-\delta(I\cap J)=0.$
\end{proof}
\begin{lem}
\label{Lemma. Irreducible BK-sets don't intersect}In a reducible
BK-polymatroid, irreducible BK-sets do not intersect.
\end{lem}

\begin{proof}
The empty polymatroid can be considered as a BK-polymatroid: the defect
of the empty set is zero, because the polymatroid rank function $f$
is normalized, $f(\varnothing)=0$.
\end{proof}
\begin{defn}
For a BK-polymatroid $P$, a \textit{BK-filtration} is a sequence
of polymatroids $P_{1}\subset P_{2}\subset...\subset P_{k}=P$ with
ground sets $E_{1}\subset E_{2}\subset...\subset E_{k}=E$ such that
all contractions $P_{i}/E_{i-1}$ are BK-polymatroids. A BK-filtration
is \textit{maximal} if all BK-polymatroids $P_{i}/E_{i-1}$ are irreducible.
\end{defn}

\begin{cor}
For a BK-polymatroid, the number of subsets in a maximal BK-filtration
equals the cardinality of the Birkhoff poset.
\end{cor}

For a BK-polymatroid $P,$ the Birkhoff poset $\mathtt{P}$ defines
a partition of the ground set $E$ as follows. Every element $\alpha$
of the poset $\mathtt{P}$ corresponds to some subset $K_{\alpha}\subseteq E$,
and every order ideal $\mathtt{I}$ of the poset $\mathtt{P}$ corresponds
to some BK-set $K_{\mathtt{I}}=\underset{\alpha\in\mathtt{I}}{\sqcup}K_{\alpha}$.
Denote by $\hat{P}_{\alpha}$ the polymatroid minor $P_{(\alpha)}/K_{(\alpha)\backslash\alpha}$,
where $P_{(\alpha)}$ is the restriction of the polymatroid $P$ on
the BK-set $K_{(\alpha)}$.
\begin{thm}
\label{Theorem. BK-partition}The ground set of a reducible BK-polymatroid
admits a unique partition into subsets $E=\underset{\alpha\in\mathtt{P}}{\sqcup}K_{\alpha}$
such that all polymatroids $\hat{P}_{\alpha}$ are irreducible BK-polymatroids
for every element $\alpha$ from the Birkhoff poset $\mathtt{P}$.
\end{thm}

\begin{proof}
Minimal elements of the poset $\mathtt{P}$ correspond to irreducible
BK-sets of the ground set $E$. By Lemma \ref{Lemma. Irreducible BK-sets don't intersect},
irreducible BK-sets do not intersect in an independent set. Every
element $\alpha$ from the poset $\mathtt{P}$ is associated with
the subset $K_{(\alpha)}\backslash K_{(\alpha)\backslash\alpha}$,
where $(\alpha)$ and $(\alpha)\backslash\alpha$ are order ideals
that correspond to BK-sets $K_{(\alpha)}$ and $K_{(\alpha)\backslash\alpha}$.
The polymatroid $\hat{P}_{\alpha}$ is the contraction by the BK-set
$K_{(\alpha)\backslash\alpha}$ for the restriction of the polymatroid
$P$ on the BK-set $K_{(\alpha)}$, and it is a BK-polymatroid by
Lemma \ref{Lemma. Contraction by BK does not change defect} for every
element $\alpha$ . The BK-polymatroid $\hat{P}_{\alpha}$ is irreducible,
because, otherwise, it would not correspond to an element in the poset
$\mathtt{P}$. The partition is unique by the construction.
\end{proof}

\section{Realizable polymatroids}

A polymatroid $P$ is \textit{realizable} over a field if there exists
a \textit{realization} via a tuple of vector subspaces $(L_{1},...,L_{n})$
in a vector space $V$ and the rank function $f(I)=dim\,L_{I}$, $L_{I}=\underset{i\in I}{\sum}L_{i}$. 

There is one more characterization of independent subsets in realizable
polymatroids via their realizations.
\begin{thm}
A tuple of vector subspaces $(L_{1},...,L_{n})$ corresponds to an
independent realizable polymatroid if and only if there exists a linearly
independent set of vectors $\upsilon_{1}\in L_{1}$, ..., $\upsilon_{n}\in L_{n}$.
\end{thm}

\begin{proof}
See Theorem 4, \cite{khovanskii_newton_2016}.
\end{proof}

\subsection{Partition of a vector space by a realizable polymatroid}

A \textit{dual realization }of a polymatroid $P$ is called a tuple
of subspaces $(L_{1},...,L_{n})$ with the rank function $f(I)=codim\,L_{I}$,
$L_{I}=\underset{i\in I}{\cap}L_{i}$. Both realizations correspond
to the same polymatroid.

In this subsection, we consider polymatroids with a lattice of flats
$\mathcal{L}$ and their dual realizations.
\begin{lem}
\label{Lemma. For a point x, I is flat}For a point $x\in V,$ the
set $I=\{i\in[n]\,|\;L_{i}\ni x\}$ is flat.
\end{lem}

\begin{proof}
If the set $I$ is not flat, there exists an element $j\in F\backslash I$
such that $F=cl_{P}(I).$ Then the point $x$ does not lie in the
subspaces $x\notin L_{j}$ and $x\notin L_{F}$. We obtain the contradiction
since $L_{I}\neq L_{F}$.
\end{proof}
\begin{lem}
\label{Lemma. B_F}The set of points in the vector space $V$ corresponding
to the same flat $F$ is
\[
B_{F}=L_{F}\backslash\underset{F'\in\mathcal{L}\backslash(F)}{\cup}L_{F'},
\]
where $(F)$ is the principal order ideal in the lattice of flats
$\mathcal{L}$.
\end{lem}

\begin{proof}
The points corresponding to a flat $F$ belong to the subspace $L_{F}$.
However, we need to remove points that do not correspond to the flat
$F$. For an inclusion of flats $F'\subset F$, there exists an inverse
inclusion of subspaces $L_{F'}\supset L_{F}$. Hence the subspace
$L_{F}$ lies in the subspace $L_{F'}$ only for flats $F'$ from
the principal order ideal $(F)$ in the lattice $\mathcal{L}$, and
it does not lie in the subspaces of other flats. This determines the
set $B_{F}=L_{F}\backslash\underset{F'\in\mathcal{L}\backslash(F)}{\cup}L_{F'}.$
\end{proof}
\begin{prop}
\textup{\label{Proposition. Polymatroid partition of a vector space}}For
a tuple of subspaces $(L_{1},...,L_{n})$ from a vector space $V$
over an infinite field, there exists a unique partition of the space
$V$ into constructible subsets $B_{F}$, enumerated by the lattice
of flats $\mathcal{L}$ of the polymatroid $P$, $V=\underset{F\in\mathcal{L}}{\sqcup}B_{F}$.
\end{prop}

\begin{proof}
Since every point of the vector space $V$ corresponds to a unique
flat, Lemma \ref{Lemma. For a point x, I is flat} provides a set-theoretic
map $\gamma:V\rightarrow\mathcal{L}$. This map defines a partition
of the space into subsets $\gamma^{-1}(F)=B_{F}$ by Lemma \ref{Lemma. B_F}.
For a finite field, the set $B_{F}$ can be empty. For an infinite
field the set $B_{F}$ is never empty as a finite intersection of
dense subsets $B_{F}=\underset{F'\in\mathcal{L}\backslash(F)}{\cap}(L_{F}\backslash L_{F'}).$
\end{proof}

\subsection{Quotients and realizable contractions}

In this subsection, we use a realization $\mathsf{n}=(L_{1},...,L_{n})$
of a polymatroid with the rank function $f(I)=dim\,\langle\mathsf{n}_{I}\rangle$,
where a subtuple $\mathsf{n}_{I}=(L_{i},\;i\in I)$ and its sum $\langle\mathsf{n}_{I}\rangle=\underset{i\in I}{\sum}L_{i}$.
\begin{defn}
For subtuples $\mathsf{n}_{I}$ and $\mathsf{n}_{J}$ from a tuple
of vector subspaces $\mathsf{n}$ in a vector space $V$, the \textit{quotient
tuple} is defined as $\mathsf{n}_{J}/\mathsf{n}_{I}=\pi(\mathsf{n}_{J\backslash I})$
for the projection $\pi:V\rightarrow V/\langle\mathsf{n}_{I}\rangle$.
\end{defn}

\begin{prop}
\label{Proposition. Polymatroid contraction and quotient}Let a tuple
of subspaces $\mathsf{n}$ be a realization of a polymatroid $P$;
then the polymatroid contraction $P/I$ for a subset $J$ by a subset
$I$ admits a realization as the quotient $\mathsf{n}_{J}/\mathsf{n}_{I}$.
\end{prop}

\begin{proof}
The polymatroid rank function for the contraction $P/I$ is equals
to the dimension of the quotient tuple of a realization $\mathsf{n}$:
\begin{align*}
rk_{I}(J) & =rk(I\cup J)-rk(I)=dim\,\langle\mathsf{n}_{I\cup J}\rangle-dim\,\langle\mathsf{n}_{I}\rangle=dim\,\langle\mathsf{n}_{I\cup J}\rangle/\langle\mathsf{n}_{I}\rangle=\\
 & =dim\,\langle\mathsf{n}_{J\backslash I}\rangle/\langle\mathsf{n}_{I}\rangle=dim\,\pi(\langle\mathsf{n}_{J\backslash I}\rangle)=dim\,\langle\pi(\mathsf{n}_{J\backslash I})\rangle=dim\,\langle\mathsf{n}_{J}/\mathsf{n}_{I}\rangle,
\end{align*}
where $I$ and $J$ are subsets from the ground set of the polymatroid
$P$.
\end{proof}
\begin{rem}
1) For a polymatroid $P$, we have a coherence between a polymatroid
realization $\mathsf{n}$, a polymatroid contraction $P/I$, and the
quotient $\mathsf{n}/\mathsf{n}_{I}$:
\[
\xymatrix{P\ar[rr]^{realization}\ar[d]_{contraction} &  & \mathsf{n}\ar[d]^{quotient}\\
P/I\ar[rr]^{realization} &  & \mathsf{n}/\mathsf{n}_{I}.
}
\]

2) We suggest the separate notion of quotient tuples of vector subspaces
to emphasize that polymatroid contractions of realizable polymatroids
can be expressed as quotients in linear algebra, in contrast to the
abstract case.
\end{rem}


\section{\label{Section. Disctriminants}Discriminants of polynomial systems}

The structure of the induced matroid on a tuple of Newton polytopes
sets up a new approach to the study of polynomial systems.

The notion of essential sets arose from the work of Sturmfels about
resultants \cite{sturmfels_newton_1994}. The idea of BK-sets arose
from the Kouchnirenko-Bernstein theorem \cite{bernshtein_number_1975}
and the Minkowski criteria about zeroing mixed volumes. The notion
of polymatroid contractions arose from the substitution of a solution
from a polynomial subsystem into the complement subsystem.

Corollary \ref{Corollary. Defect of circuits is -1} links mixed discriminants
\cite{cattani_mixed_2013} and induced matroids with a unique circuit.
Theorem \ref{Theorem. Maximal essential =00003D minimal defect} provides
a new characterization of the tuple of Newton polytopes for resultants
\cite{sturmfels_newton_1994}. In a generic case, the defect for the
basis of the induced matroid (Theorem \ref{Theorem. Basis have the same defect})
equals the expected dimension of the set of solutions for the system
in the non-empty case.

For a BK-tuple of Newton polytopes, the Birkhoff poset provides the
optimal algorithm for solving the polynomial systems. First, we solve
subsystems corresponding to the minimal elements of the poset. Then
substitute solutions into the remaining system and repeat the procedure.
This algorithm already appeared in \cite{brysiewicz_solving_2021}
and the Macaulay2 package $\mathtt{DecomposableSparseSystems}$.

The polymatroid partition of a vector space in Proposition \ref{Proposition. Polymatroid partition of a vector space}
is motivated by the Esterov conjecture about irreducibility of discriminants
\cite{esterov_galois_2019,pokidkin_components_2025}. The conjecture
is crucial for a general description of discriminants for BK-tuples
of Newton polytopes \cite{pokidkin_components_2025}. The combinatorial
part of the description is encoded by Theorem \ref{Theorem. BK-partition}.

\subsubsection*{Acknowledgment}

I am grateful to Alexander Esterov for supervision and to my family
for support.

\subsubsection*{Funding}

The study was partially supported by the HSE University Basic Research
Program.

\begin{adjustwidth}{-0.0in}{-0.0in}
	\sloppy
	\printbibliography
	\fussy
\end{adjustwidth}

National Research University Higher School of Economics, Russian Federation

AG Laboratory, HSE, 6 Usacheva str., Moscow, Russia, 119048

\-

Email: vppokidkin@hse.ru
\end{document}
